\documentclass[11pt,a4paper]{article}

% --- Packages ---
\usepackage[utf8]{inputenc}
\usepackage[T1]{fontenc}
\usepackage[spanish]{babel}
\usepackage{geometry}
\geometry{margin=2.5cm}
\usepackage{hyperref}
\hypersetup{colorlinks=true, linkcolor=blue, urlcolor=blue}
\usepackage{booktabs}
\usepackage{longtable}
\usepackage{array}
\usepackage{fancyhdr}
\usepackage{titlesec}
\usepackage{enumitem}
\usepackage{xcolor}
\usepackage{tabularx}
\usepackage{graphicx}
\usepackage{multicol}
\usepackage{float}
\usepackage{amsmath}

% --- Colors ---
\definecolor{villafox}{HTML}{D97706}
\definecolor{villablack}{HTML}{0A0A0A}
\definecolor{sectioncolor}{HTML}{1a1a1a}

% --- Header/Footer ---
\pagestyle{fancy}
\fancyhf{}
\fancyhead[L]{\small\textcolor{gray}{PRD -- La Villa Skateboarding API}}
\fancyhead[R]{\small\textcolor{gray}{Confidencial}}
\fancyfoot[C]{\thepage}
\renewcommand{\headrulewidth}{0.4pt}

% --- Section formatting ---
\titleformat{\section}{\Large\bfseries\color{sectioncolor}}{}{0em}{}[\titlerule]
\titleformat{\subsection}{\large\bfseries\color{sectioncolor}}{}{0em}{}
\titleformat{\subsubsection}{\normalsize\bfseries\color{sectioncolor}}{}{0em}{}

% --- Custom commands ---
\newcommand{\apiversion}{v1.0}
\newcommand{\docdate}{01 de Septiembre de 2026}
\newcommand{\contractvalue}{\$1.000 a \$1.500\ USD (estimado)}
\newcommand{\contractperiod}{12 meses}

\begin{document}

% ========================================
% COVER PAGE
% ========================================
\begin{titlepage}
\centering
\vspace*{2cm}

{\Huge\bfseries\textcolor{sectioncolor}{PRD}\\[0.3cm]}
{\Large\textcolor{sectioncolor}{Product Requirements Document}\\[1.5cm]}

{\Large\textcolor{villafox}{\textbf{LA VILLA SKATEBOARDING}}}\\[0.5cm]
{\Large Sistema de API Microservicios para E-Commerce}\\[2cm]

\rule{\textwidth}{1.5pt}\\[1cm]

\begin{minipage}{0.45\textwidth}
\begin{flushleft}
\textbf{\textcolor{villafox}{Cliente:}}\\
La Villa Skateboarding\\
Nariño, Colombia\\[1cm]
\textbf{\textcolor{villafox}{Versión del Documento:}}\\
\apiversion\\
\end{flushleft}
\end{minipage}
\hfill
\begin{minipage}{0.45\textwidth}
\begin{flushright}
\textbf{\textcolor{villafox}{Proveedor:}}\\
MACREAT - developer\\[1cm]
\textbf{\textcolor{villafox}{Fecha:}}\\
\docdate\\
\end{flushright}
\end{minipage}

\vfill

\begin{minipage}{0.8\textwidth}
\centering
\textcolor{gray}{\textit{Este documento es confidencial y está destinado únicamente a las partes involucradas en el proyecto.}}
\end{minipage}

\vspace{1cm}
\end{titlepage}

% ========================================
% TABLE OF CONTENTS
% ========================================
\tableofcontents
\newpage

% ========================================
% 1. INFORMACIÓN DEL CONTRATO
% ========================================
\section{Información del Contrato}

\begin{tabularx}{\textwidth}{lX}
\toprule
\textbf{Campo} & \textbf{Detalle} \\
\midrule
Objeto & Diseño, construcción, levantamiento y mantenimiento de API microservicios para e-commerce \\
Valor Total & \contractvalue \\
Duración & \contractperiod \\
Moneda & Dólares Americanos (USD) \\
Fecha de Inicio & Por definir al momento de la firma \\
Entorno & Producción: VPS + Vercel \\
\bottomrule
\end{tabularx}

% ========================================
% 2. ALCANCE DEL PROYECTO
% ========================================
\section{Alcance del Proyecto}

\subsection{Descripción General}

El sistema consiste en una arquitectura de microservicios para una tienda en línea de skateboarding, compuesta por:

\begin{itemize}[leftmargin=2cm]
    \item \textbf{API Gateway} (Laravel/PHP) -- Enrutamiento, autenticación, proxy a microservicios
    \item \textbf{6 Microservicios} (Python/FastAPI) -- Catálogo, Inventario, Órdenes, Pagos, Notificaciones, Usuarios
    \item \textbf{Frontend} (Next.js/React) -- Tienda pública + panel administrativo
    \item \textbf{Infraestructura} (Docker Compose) -- PostgreSQL, Redis, MinIO, RabbitMQ, MeiliSearch
\end{itemize}

\subsection{Arquitectura}

\begin{verbatim}
Frontend (Next.js :3000)
        |
        v
  Gateway (Laravel :8010) ---[Sanctum Auth]---> PostgreSQL
        |
        +---> Catalog (:9002)    ---> PostgreSQL + MinIO
        +---> Inventory (:9003)  ---> PostgreSQL
        +---> Orders (:9004)     ---> PostgreSQL + Redis
        +---> Payments (:9005)   ---> PostgreSQL
        +---> Notifications (:9006) -> WhatsApp API + Twilio SMS
        +---> Users (:9007)      ---> PostgreSQL
\end{verbatim}

% ========================================
% 3. SERVICIOS Y ENDPOINTS
% ========================================
\section{Especificación de Servicios y Endpoints}

\subsection{API Gateway (Laravel)}

\begin{tabularx}{\textwidth}{|l|X|l|}
\hline
\textbf{Método} & \textbf{Ruta} & \textbf{Auth} \\
\hline
POST & \texttt{/api/admin/login} & No \\
POST & \texttt{/api/admin/logout} & Sanctum \\
GET & \texttt{/api/admin/me} & Sanctum \\
PUT & \texttt{/api/admin/me} & Sanctum \\
\hline
GET & \texttt{/api/v1/catalog/\{path\}} & No \\
POST & \texttt{/api/v1/catalog/search} & No \\
GET & \texttt{/api/v1/\{service\}/health} & No \\
POST & \texttt{/api/v1/\{service\}/\{path\}} & No \\
ANY & \texttt{/api/v1/\{service\}/\{path\}} & Sanctum \\
\hline
\end{tabularx}

\subsubsection{Autenticación}
\begin{itemize}
    \item Laravel Sanctum con tokens scoped \texttt{['admin']}
    \item Middleware \texttt{EnsureAdmin}: verifica \texttt{user.is\_admin} y \texttt{tokenCan('admin')}
    \item Login: email + password + device\_name, retorna \texttt{plainTextToken}
    \item Rate limit: 5 intentos/minuto en login
\end{itemize}

\subsection{Servicio de Catálogo (FastAPI :9002)}

\begin{tabularx}{\textwidth}{|l|X|l|}
\hline
\textbf{Método} & \textbf{Ruta} & \textbf{Descripción} \\
\hline
GET & \texttt{/health} & Health check \\
GET & \texttt{/products} & Listar productos (parámetros: search, id) \\
GET & \texttt{/products/\{id\}} & Obtener producto por ID \\
POST & \texttt{/products} & Crear producto \\
POST & \texttt{/search} & Búsqueda inteligente de productos \\
POST & \texttt{/media/presign} & URL presignada para upload S3 \\
POST & \texttt{/media/\{id\}/complete} & Marcar upload como completado \\
POST & \texttt{/subscribers} & Suscribir email \\
GET & \texttt{/subscribers} & Listar suscriptores \\
GET & \texttt{/subscribers/count} & Conteo de suscriptores \\
\hline
\end{tabularx}

\subsubsection{Estrategia de Búsqueda Inteligente}
La búsqueda inteligente implementa una estrategia de filtrado por palabras clave con las siguientes capas:

\begin{enumerate}
    \item \textbf{Extracción de palabras clave}: El texto libre se tokeniza y se filtran palabras vacías (español e inglés: ``de'', ``el'', ``la'', ``para'', ``the'', ``is'', etc.)
    \item \textbf{Stemming español básico}: Se eliminan terminaciones plurales para mejorar recall (``negros'' $\rightarrow$ ``negro'', ``ruedas'' $\rightarrow$ ``rueda''). Se aplica solo a palabras con 5+ caracteres para evitar distorsión de términos cortos como ``tenis''
    \item \textbf{Lógica OR}: Un producto coincide si contiene CUALQUIERA de las palabras clave extraídas (no todas)
    \item \textbf{Scoring de relevancia}: Cada coincidencia suma 1 punto. Los resultados se ordenan por puntaje descendente, luego por nombre
    \item \textbf{Campos de búsqueda}: nombre del producto, nombre de categoría, SKU, descripción
    \item \textbf{Filtros duros}: categoría, subcategoría, rango de precio, marca (cuando se proporcionan como filtros estructurados)
\end{enumerate}

\textbf{Ejemplo}: La consulta ``tenis nike negros'' extrae las keywords [tenis, nike, negros], aplica stemming [tenis, nike, negro], y retorna todos los productos donde CUALQUIERA de esas palabras aparece en nombre/categoría/descripción, ordenados por relevancia.

\subsubsection{Tablas de Base de Datos}
\begin{itemize}
    \item \texttt{categories}: id, slug (unique), name, category\_group
    \item \texttt{products}: id, name, description, sku (unique), price (Decimal 10,2), category\_id (FK), media\_id (FK), active
    \item \texttt{media}: id, storage\_key (unique), bucket, mime\_type, status, variants (JSON), metadata (JSON)
    \item \texttt{subscribers}: id, email (unique), tag
\end{itemize}

\subsection{Servicio de Inventario (FastAPI :9003)}

\begin{tabularx}{\textwidth}{|l|X|l|}
\hline
\textbf{Método} & \textbf{Ruta} & \textbf{Descripción} \\
\hline
GET & \texttt{/health} & Health check \\
GET & \texttt{/inventory} & Listar niveles de stock \\
GET & \texttt{/inventory/\{product\_id\}} & Stock de un producto \\
PUT & \texttt{/inventory/\{product\_id\}} & Establecer cantidad \\
POST & \texttt{/inventory/\{product\_id\}/adjust} & Ajustar stock por delta \\
GET & \texttt{/availability} & Verificar disponibilidad (IDs separados por coma) \\
\hline
\end{tabularx}

\subsubsection{Tablas}
\begin{itemize}
    \item \texttt{stock\_levels}: id, product\_id (unique), sku, quantity, low\_stock\_threshold (default: 5)
\end{itemize}

\subsection{Servicio de Órdenes (FastAPI :9004)}

\begin{tabularx}{\textwidth}{|l|X|l|}
\hline
\textbf{Método} & \textbf{Ruta} & \textbf{Descripción} \\
\hline
GET & \texttt{/health} & Health check \\
POST & \texttt{/orders} & Crear orden directa \\
GET & \texttt{/orders} & Listar órdenes (skip/limit) \\
GET & \texttt{/orders/summary} & Resumen dashboard \\
GET & \texttt{/orders/\{id\}} & Obtener orden \\
POST & \texttt{/orders/\{id\}/cancel} & Cancelar orden \\
POST & \texttt{/orders/\{id\}/status} & Actualizar estado \\
\hline
GET & \texttt{/cart} & Obtener carrito (requiere \texttt{X-Session-Id}) \\
POST & \texttt{/cart/items} & Agregar item al carrito \\
DELETE & \texttt{/cart/items/\{product\_id\}} & Eliminar item \\
POST & \texttt{/cart/checkout} & Checkout: crear orden desde carrito \\
\hline
\end{tabularx}

\subsubsection{Estados de Orden}
\texttt{pending} $\rightarrow$ \texttt{confirmed} $\rightarrow$ \texttt{shipped} $\rightarrow$ \texttt{delivered} $\mid$ \texttt{cancelled}

\subsubsection{Carrito}
\begin{itemize}
    \item Almacenamiento: Redis con TTL configurable
    \item Clave: \texttt{cart:\{session\_id\}}
    \item Checkout: convierte carrito a Orden persistente, elimina carrito
\end{itemize}

\subsection{Servicio de Pagos (FastAPI :9005)}

\begin{tabularx}{\textwidth}{|l|X|l|}
\hline
\textbf{Método} & \textbf{Ruta} & \textbf{Descripción} \\
\hline
GET & \texttt{/health} & Health check \\
POST & \texttt{/payment-intents} & Crear intención de pago \\
GET & \texttt{/payment-intents} & Listar intenciones \\
GET & \texttt{/payment-intents/\{id\}} & Obtener intención \\
POST & \texttt{/payment-intents/\{id\}/confirm} & Confirmar pago \\
POST & \texttt{/payment-intents/\{id\}/capture} & Capturar pago \\
\hline
\end{tabularx}

\subsubsection{Tablas}
\begin{itemize}
    \item \texttt{payment\_intents}: id, order\_id, amount (float), currency (default: ``COP''), status (pending/confirmed/captured), provider (default: ``manual'')
\end{itemize}

\subsection{Servicio de Notificaciones (FastAPI :9006)}

\begin{tabularx}{\textwidth}{|l|X|l|}
\hline
\textbf{Método} & \textbf{Ruta} & \textbf{Descripción} \\
\hline
GET & \texttt{/health} & Health check \\
POST & \texttt{/send-catalog} & Enviar catálogo completo vía WhatsApp \\
POST & \texttt{/send-whatsapp} & Enviar mensaje WhatsApp arbitrario \\
POST & \texttt{/notify-order} & Notificar orden vía SMS (Twilio) \\
GET & \texttt{/status} & Estado de configuración WhatsApp/Twilio \\
\hline
\end{tabularx}

\subsubsection{Integraciones Externas}
\begin{itemize}
    \item WhatsApp Cloud API (Meta) v21.0
    \item Twilio SMS (fallback)
    \item Modo mock: responde éxito simulado si no hay API keys configuradas
\end{itemize}

\subsection{Servicio de Usuarios (FastAPI :9007)}

\begin{tabularx}{\textwidth}{|l|X|l|}
\hline
\textbf{Método} & \textbf{Ruta} & \textbf{Descripción} \\
\hline
GET & \texttt{/health} & Health check \\
GET & \texttt{/users} & Listar usuarios \\
GET & \texttt{/users/\{id\}} & Obtener usuario \\
POST & \texttt{/users} & Crear usuario (email único) \\
PUT & \texttt{/users/\{id\}} & Actualizar usuario \\
\hline
\end{tabularx}

% ========================================
% 4. INFRAESTRUCTURA
% ========================================
\section{Infraestructura}

\subsection{Servicios Docker}

\begin{tabularx}{\textwidth}{|l|l|l|X|}
\hline
\textbf{Servicio} & \textbf{Imagen} & \textbf{Puerto} & \textbf{Propósito} \\
\hline
PostgreSQL & postgres:16-alpine & 5432 & Base de datos principal \\
Redis & redis:7-alpine & 6379 & Carrito, caché \\
RabbitMQ & rabbitmq:3-mgmt & 5672 & Cola de mensajes \\
MeiliSearch & meilisearch:v1 & 7700 & Motor de búsqueda \\
MinIO & minio/minio & 9000/9001 & Almacenamiento S3 \\
Gateway & Laravel & 8010 & API Gateway \\
Catalog & FastAPI & 9002 & Catálogo de productos \\
Inventory & FastAPI & 9003 & Gestión de stock \\
Orders & FastAPI & 9004 & Órdenes + carrito \\
Payments & FastAPI & 9005 & Intenciones de pago \\
Notifications & FastAPI & 9006 & WhatsApp + SMS \\
Users & FastAPI & 9007 & Perfiles de usuario \\
Frontend & Next.js & 3000 & Tienda + admin \\
\hline
\end{tabularx}

\subsection{Volúmenes Persistidos}
\begin{itemize}
    \item \texttt{postgres\_data}
    \item \texttt{redis\_data}
    \item \texttt{rabbitmq\_data}
    \item \texttt{meilisearch\_data}
    \item \texttt{minio\_data}
\end{itemize}

\subsection{Variables de Entorno Críticas}

\begin{tabularx}{\textwidth}{|l|X|}
\hline
\textbf{Variable} & \textbf{Descripción} \\
\hline
\texttt{DATABASE\_URL} & URL de conexión PostgreSQL \\
\texttt{NEXT\_PUBLIC\_API\_BASE\_URL} & URL pública del gateway \\
\texttt{NEXT\_PUBLIC\_SITE\_URL} & Dominio del sitio \\
\texttt{FRONTEND\_URL} & Dominio para CORS \\
\texttt{SANCTUM\_STATEFUL\_DOMAINS} & Dominios Sanctum \\
\texttt{ADMIN\_PASSWORD} & Contraseña del admin \\
\texttt{APP\_KEY} & Clave de cifrado Laravel \\
\texttt{CDN\_BASE\_URL} & URL pública de MinIO/S3 \\
\hline
\end{tabularx}

% ========================================
% 5. FUNCIONALIDADES CLAVE
% ========================================
\section{Funcionalidades Clave}

\subsection{Tienda Pública}
\begin{enumerate}
    \item Listado de productos con filtrado por categoría y subcategoría
    \item Búsqueda inteligente de productos por palabras clave con stemming
    \item Carrito de compras persistente (sesión + Redis)
    \item Checkout con creación de orden
    \item Detalle de producto con imágenes
    \item Sistema de suscripción por email
    \item Contenido de marca: hero, editorial, heritage, identidad visual
\end{enumerate}

\subsection{Panel Administrativo}
\begin{enumerate}
    \item Dashboard con métricas (órdenes totales, ingresos, hoy/ayer)
    \item Gestión de productos (CRUD + imágenes via presigned URLs)
    \item Gestión de órdenes (estados, cancelación)
    \item Gestión de inventario (stock, ajustes, alertas de bajo stock)
    \item Autenticación con tokens Sanctum
\end{enumerate}

\subsection{Integraciones Externas}
\begin{enumerate}
    \item WhatsApp Cloud API (Meta) para notificaciones y catálogo
    \item Twilio SMS para notificaciones de órdenes
    \item MinIO (S3) para almacenamiento de imágenes
\end{enumerate}

\subsection{Búsqueda Inteligente (Actual)}
\begin{itemize}
    \item Estrategia basada en keywords con stemming español
    \item Categorías: decks, apparel, accessories, gear
    \item Subcategorías: 14 tipos (tenis, buzos, camisetas, trucks, ruedas, etc.)
    \item Extracción de filtros: categoría, subcategoría, precio, marca
    \item Lógica OR con scoring de relevancia por coincidencia
    \item Sin dependencia de servicios externos para búsqueda
\end{itemize}

% ========================================
% 6. SEGURIDAD
% ========================================
\section{Seguridad}

\begin{itemize}
    \item Autenticación basada en tokens (Laravel Sanctum)
    \item Tokens scoped con ability \texttt{['admin']}
    \item Middleware \texttt{EnsureAdmin} en rutas protegidas
    \item Rate limiting: login (5/min), búsqueda NL (10/min), checkout (20/min)
    \item Variables de entorno para credenciales (no hardcodeadas en código fuente)
    \item CORS configurado por dominio
    \item HTTPS requerido en producción
\end{itemize}

% ========================================
% 7. ENTREGABLES
% ========================================
\section{Entregables}

\subsection{Fase 1: Diseño (Semanas 1--2)}
\begin{itemize}
    \item Documento de requerimientos (este documento)
    \item Arquitectura de microservicios
    \item Diseño de base de datos
    \item Diagrama de endpoints
    \item Definición de variables de entorno
\end{itemize}

\subsection{Fase 2: Construcción (Semanas 3--8)}
\begin{itemize}
    \item API Gateway (Laravel) con autenticación
    \item 6 microservicios FastAPI
    \item Frontend Next.js (tienda + admin)
    \item Infraestructura Docker Compose
    \item Integración WhatsApp/SMS
    \item Búsqueda inteligente de productos
\end{itemize}

\subsection{Fase 3: Levantamiento (Semanas 9--10)}
\begin{itemize}
    \item Configuración de VPS (Docker, nginx, TLS)
    \item Despliegue frontend en Vercel
    \item Configuración de variables de entorno producción
    \item Generación de APP\_KEY único
    \item Pruebas end-to-end
\end{itemize}

\subsection{Fase 4: Mantenimiento (Mes 1--12)}
\begin{itemize}
    \item Monitoreo de servicios
    \item Resolución de bugs
    \item Actualizaciones de seguridad
    \item Soporte técnico
    \item Mejoras menores (20--30 horas/mes)
\end{itemize}

% ========================================
% 8. CRONOGRAMA
% ========================================
\section{Cronograma}

\begin{tabularx}{\textwidth}{|l|l|X|}
\hline
\textbf{Fase} & \textbf{Duración} & \textbf{Actividades} \\
\hline
Diseño & 2 semanas & PRD, arquitectura, DB, endpoints \\
Construcción & 6 semanas & Gateway, microservicios, frontend \\
Levantamiento & 2 semanas & VPS, Vercel, TLS, pruebas \\
Mantenimiento & 12 meses & Monitoreo, bugs, seguridad, soporte \\
\hline
\end{tabularx}

% ========================================
% 9. TERMINOS ECONÓMICOS
% ========================================
\section{Terminos Económicos}

\begin{tabularx}{\textwidth}{|l|X|}
\hline
\textbf{Concepto} & \textbf{Detalle} \\
\hline
Valor total del contrato & \contractvalue \\
Forma de pago & 50\% al inicio, 50\% al entrega \\
Moneda & Dólares Americanos (USD) \\
Periodo de mantenimiento & \contractperiod \\
Incluye mantenimiento & Sí (20--30 horas/mes) \\
\hline
\end{tabularx}

% ========================================
% 10. total de endpoints
% ========================================
\section{Resumen de Endpoints}

\begin{tabularx}{\textwidth}{|l|c|X|}
\hline
\textbf{Servicio} & \textbf{Endpoints} & \textbf{Descripción} \\
\hline
Gateway & 9 & Enrutamiento, auth, proxy \\
Catálogo & 10 & Productos, búsqueda NL, media, suscriptores \\
Inventario & 6 & Stock, disponibilidad \\
Órdenes & 10 & Órdenes, carrito, checkout \\
Pagos & 5 & Intenciones de pago \\
Notificaciones & 5 & WhatsApp, SMS \\
Usuarios & 5 & CRUD de usuarios \\
\hline
\textbf{Total} & \textbf{50} & \\
\hline
\end{tabularx}

% ========================================
% 11. DIAGRAMAS
% ========================================
\section{Diagramas}

\subsection{Diagrama de Arquitectura Lógica}

Muestra la estructura del sistema por capas: presentación, gateway, microservicios y datos.

\begin{verbatim}
+====================================================================+
|                      PRESENTACION (Capa de UI)                     |
|  +-------------------+        +--------------------------------+  |
|  | Tienda Publica    |        | Panel Administrativo           |  |
|  | (Next.js :3000)   |        | (Next.js :3000/admin)          |  |
|  +-------------------+        +--------------------------------+  |
+====================================================================+
                            |  HTTP/HTTPS
                            v
+====================================================================+
|                    GATEWAY (Capa de Enrutamiento)                  |
|  +--------------------------------------------------------------+  |
|  | API Gateway (Laravel :8010)                                  |  |
|  | - Autenticacion Sanctum                                      |  |
|  | - Rate Limiting                                              |  |
|  | - CORS / Headers                                             |  |
|  | - Proxy a microservicios                                     |  |
|  +--------------------------------------------------------------+  |
+====================================================================+
              |              |              |              |
              v              v              v              v
+====================================================================+
|                 MICROSERVICIOS (Capa de Logica)                    |
|  +-----------+ +-----------+ +-----------+ +-----------+           |
|  | Catalog   | | Inventory | | Orders    | | Payments  |           |
|  | :9002     | | :9003     | | :9004     | | :9005     |           |
|  +-----------+ +-----------+ +-----------+ +-----------+           |
|  +-----------+ +-----------+                                      |
|  | Notifs    | | Users     |                                      |
|  | :9006     | | :9007     |                                      |
|  +-----------+ +-----------+                                      |
+====================================================================+
              |              |              |
              v              v              v
+====================================================================+
|                      DATOS (Capa de Persistencia)                  |
|  +--------------+  +--------------+  +--------------+              |
|  | PostgreSQL   |  | Redis        |  | MinIO (S3)   |              |
|  | :5432        |  | :6379        |  | :9000        |              |
|  | - products   |  | - carrito    |  | - imagenes   |              |
|  | - orders     |  | - cache      |  | - media      |              |
|  | - users      |  | - sessions   |  |              |              |
|  +--------------+  +--------------+  +--------------+              |
+====================================================================+
\end{verbatim}

\subsection{Diagrama de Arquitectura Funcional}

Muestra las funcionalidades agrupadas por dominio de negocio.

\begin{verbatim}
+====================================================================+
|                     DOMINIO: TIENDA PUBLICA                        |
|  +------------------+  +------------------+  +------------------+  |
|  | Catalogo         |  | Busqueda         |  | Carrito           |  |
|  | - Listar prod.   |  | - Por keywords   |  | - Agregar item   |  |
|  | - Ver detalle    |  | - Por categoria  |  | - Eliminar item  |  |
|  | - Filtrar        |  | - Por precio     |  | - Ver carrito    |  |
|  +------------------+  +------------------+  +------------------+  |
|  +------------------+  +------------------+  +------------------+  |
|  | Checkout         |  | Suscripcion      |  | Contenido Marca  |  |
|  | - Crear orden    |  | - Email signup   |  | - Hero           |  |
|  | - Pago           |  | - Newsletter     |  | - Editorial      |  |
|  +------------------+  +------------------+  +------------------+  |
+====================================================================+

+====================================================================+
|                     DOMINIO: ADMINISTRACION                        |
|  +------------------+  +------------------+  +------------------+  |
|  | Dashboard        |  | Gestion Prod.    |  | Gestion Ordenes  |  |
|  | - Metricas       |  | - CRUD           |  | - Estados        |  |
|  | - Ingresos       |  | - Imagenes S3    |  | - Cancelar       |  |
|  | - Hoy/Ayer       |  | - Categorias     |  | - Historial      |  |
|  +------------------+  +------------------+  +------------------+  |
|  +------------------+  +------------------+                       |
|  | Inventario       |  | Auth             |                       |
|  | - Stock          |  | - Login          |                       |
|  | - Ajustes        |  | - Tokens         |                       |
|  | - Alertas        |  | - Permisos       |                       |
|  +------------------+  +------------------+                       |
+====================================================================+

+====================================================================+
|                     DOMINIO: EXTERNO                               |
|  +------------------+  +------------------+  +------------------+  |
|  | WhatsApp API     |  | Twilio SMS       |  | MinIO / S3       |  |
|  | - Notificaciones |  | - Ordenes        |  | - Upload imgs    |  |
|  | - Catalogo       |  | - Confirmaciones |  | - CDN publico    |  |
|  +------------------+  +------------------+  +------------------+  |
+====================================================================+
\end{verbatim}

\subsection{Diagrama de Endpoints por Servicio}

Mapa visual de todos los endpoints agrupados por servicio y metodo HTTP.

\begin{verbatim}
GATEWAY (Laravel :8010)
+--------------------------------------------------------------+
| [POST]   /api/admin/login              -> Login              |
| [POST]   /api/admin/logout             -> Logout (auth)      |
| [GET]    /api/admin/me                 -> Profile (auth)     |
| [PUT]    /api/admin/me                 -> Update (auth)      |
| [GET]    /api/v1/catalog/{path}        -> Proxy publico      |
| [POST]   /api/v1/catalog/search        -> Busqueda NL        |
| [GET]    /api/v1/{service}/health      -> Health check       |
| [POST]   /api/v1/{service}/{path}      -> Proxy (auth)       |
| [ANY]    /api/v1/{service}/{path}      -> Proxy (auth)       |
+--------------------------------------------------------------+

CATALOG (FastAPI :9002)
+--------------------------------------------------------------+
| [GET]    /health                        -> Health check       |
| [GET]    /products                      -> Listar productos   |
| [GET]    /products/{id}                 -> Ver producto       |
| [POST]   /products                      -> Crear producto     |
| [POST]   /search                        -> Busqueda NL        |
| [POST]   /media/presign                 -> URL upload S3      |
| [POST]   /media/{id}/complete           -> Marcar completado  |
| [POST]   /subscribers                   -> Suscribir email    |
| [GET]    /subscribers                   -> Listar suscript.   |
| [GET]    /subscribers/count             -> Conteo             |
+--------------------------------------------------------------+

INVENTORY (FastAPI :9003)
+--------------------------------------------------------------+
| [GET]    /health                        -> Health check       |
| [GET]    /inventory                     -> Listar stock       |
| [GET]    /inventory/{product_id}        -> Stock producto     |
| [PUT]    /inventory/{product_id}        -> Set cantidad       |
| [POST]   /inventory/{product_id}/adjust -> Ajustar delta      |
| [GET]    /availability                  -> Disponibilidad     |
+--------------------------------------------------------------+

ORDERS (FastAPI :9004)
+--------------------------------------------------------------+
| [GET]    /health                        -> Health check       |
| [POST]   /orders                        -> Crear orden        |
| [GET]    /orders                        -> Listar ordenes     |
| [GET]    /orders/summary                -> Dashboard          |
| [GET]    /orders/{id}                   -> Ver orden          |
| [POST]   /orders/{id}/cancel            -> Cancelar           |
| [POST]   /orders/{id}/status            -> Actualizar estado  |
| [GET]    /cart                          -> Ver carrito        |
| [POST]   /cart/items                    -> Agregar item       |
| [DELETE] /cart/items/{product_id}       -> Eliminar item      |
| [POST]   /cart/checkout                 -> Checkout           |
+--------------------------------------------------------------+

PAYMENTS (FastAPI :9005)
+--------------------------------------------------------------+
| [GET]    /health                        -> Health check       |
| [POST]   /payment-intents               -> Crear intento      |
| [GET]    /payment-intents               -> Listar             |
| [GET]    /payment-intents/{id}          -> Ver intento        |
| [POST]   /payment-intents/{id}/confirm  -> Confirmar          |
| [POST]   /payment-intents/{id}/capture  -> Capturar           |
+--------------------------------------------------------------+

NOTIFICATIONS (FastAPI :9006)
+--------------------------------------------------------------+
| [GET]    /health                        -> Health check       |
| [POST]   /send-catalog                  -> WhatsApp catalogo  |
| [POST]   /send-whatsapp                 -> WhatsApp msg       |
| [POST]   /notify-order                  -> SMS Twilio         |
| [GET]    /status                        -> Estado config      |
+--------------------------------------------------------------+

USERS (FastAPI :9007)
+--------------------------------------------------------------+
| [GET]    /health                        -> Health check       |
| [GET]    /users                         -> Listar usuarios    |
| [GET]    /users/{id}                    -> Ver usuario        |
| [POST]   /users                         -> Crear usuario      |
| [PUT]    /users/{id}                    -> Actualizar         |
+--------------------------------------------------------------+
\end{verbatim}

% ========================================
% 13. HOJA DE RUTA FUTURA
% ========================================
\section{Hoja de Ruta -- Funcionalidades Futuras}

Las siguientes funcionalidades están planificadas como mejoras independientes del alcance actual del contrato. Se implementarán bajo presupuesto adicional o en futuros ciclos de desarrollo.

\subsection{Agente Conversacional con LLM}

\begin{itemize}
    \item \textbf{Objetivo}: Asistente virtual que responda preguntas sobre la tienda (productos, disponibilidad, tallas, precios, pedidos)
    \item \textbf{Tecnología}: Modelo de lenguaje (GPT, Claude, u open-source) integrado vía API
    \item \textbf{Alcance}: Consultas sobre catálogo, estado de órdenes, recomendaciones de productos
    \item \textbf{Implementación}: Microservicio independiente (\texttt{/chat}) que consume el catálogo y devuelve respuestas contextuales
    \item \textbf{Costo}: Requiere API key del proveedor LLM (costo operativo mensual variable)
    \item \textbf{Estado}: No incluido en el contrato actual. Funcionalidad futura independiente
\end{itemize}

\subsection{Importación de Productos desde Google Drive}

\begin{itemize}
    \item \textbf{Objetivo}: Permitir al cliente subir un archivo CSV/Excel a Google Drive y que los productos se importen automáticamente al inventario
    \item \textbf{Flujo}: Google Drive API $\rightarrow$ webhook/polling $\rightarrow$ parser CSV $\rightarrow$ endpoint \texttt{/products} $\rightarrow$ registro en inventario
    \item \textbf{Requisitos}: Cuenta de servicio Google, permisos de lectura en Drive, estructura CSV predefinida
    \item \textbf{Estado}: No incluido en el contrato actual. Funcionalidad futura independiente
\end{itemize}

\subsection{Creación Directa de Productos vía API}

\begin{itemize}
    \item \textbf{Objetivo}: Endpoint público o autenticado para crear productos directamente sin pasar por el panel admin
    \item \textbf{Endpoint}: \texttt{POST /api/v1/catalog/products} (ya existe funcionalmente, requiere autenticación)
    \item \textbf{Uso}: Integración con sistemas externos, imports masivos, automatización
    \item \textbf{Estado}: Endpoint funcional. Requiere documentación y validación de schema para uso externo
\end{itemize}

% ========================================
% 14. CICLO DE DESARROLLO DE SOFTWARE
% ========================================
\section{Ciclo de Desarrollo de Software -- Modelo de Sommerville}

El desarrollo y mantenimiento de este sistema sigue el modelo de ciclo de vida de software propuesto por \textbf{Ian Sommerville} en \textit{Engineering Software} (10th ed.), que enfatiza la naturaleza \textbf{iterativa e incremental} del desarrollo moderno.

\subsection{Por Qué un Modelo Iterativo}

\begin{itemize}
    \item Los requisitos de e-commerce evolucionan constantemente (nuevas categorías, integraciones de pago, funciones de marketing)
    \item La retroalimentación del usuario es esencial para priorizar mejoras
    \item Los microservicios permiten desplegar partes del sistema de forma independiente
    \item Un modelo en cascada tradicional no se adapta a los cambios continuos del negocio
\end{itemize}

\subsection{Fases del Modelo de Sommerville}

\begin{enumerate}
    \item \textbf{Definición de Requisitos}: Se capturan y documentan los requisitos funcionales y no funcionales (este PRD). Se validan con el cliente antes de avanzar.

    \item \textbf{Diseño del Sistema}: Se define la arquitectura (microservicios, base de datos, protocolos). Se producen diagramas de componentes, secuencia y despliegue.

    \item \textbf{Implementación}: Se construyen los módulos de forma incremental. Cada microservicio se desarrolla, prueba y despliega de forma independiente.

    \item \textbf{Pruebas}: Se ejecutan pruebas unitarias, de integración y end-to-end. Se validan los criterios de aceptación de cada funcionalidad.

    \item \textbf{Despliegue}: Se libera incrementalmente. Primero se despliega en staging, luego en producción tras validación.

    \item \textbf{Evolución}: El sistema se mantiene vivo. Se corrigen bugs, se agregan funcionalidades, se actualizan dependencias. Cada ciclo de mantenimiento repite las fases anteriores de forma reducida.
\end{enumerate}

\subsection{Aplicación en Este Proyecto}

\begin{tabularx}{\textwidth}{|l|X|}
\hline
\textbf{Principio Sommerville} & \textbf{Aplicación en La Villa} \\
\hline
Desarrollo incremental & Cada microservicio se despliega y valida por separado \\
Refactorización continua & Se reestructura código sin cambiar funcionalidad \\
Pruebas automatizadas & Health checks, tests de integración, validación de endpoints \\
Retroalimentación del usuario & Iteraciones basadas en uso real de la tienda \\
Gestión de configuración & Variables de entorno, Docker Compose, despliegue reproducible \\
Documentación viva & Este PRD se actualiza con cada ciclo de desarrollo \\
\hline
\end{tabularx}

\subsection{Ciclos de Mantenimiento}

Durante los 12 meses de contrato, el mantenimiento sigue ciclos de 2--4 semanas:

\begin{enumerate}
    \item \textbf{Sprint de mantenimiento}: Se identifican bugs, mejoras menores y actualizaciones de seguridad
    \item \textbf{Desarrollo}: Se implementan los cambios en un branch dedicado
    \item \textbf{Pruebas}: Se validan los cambios contra el sistema existente
    \item \textbf{Despliegue}: Se liberan los cambios en producción
    \item \textbf{Revisión}: Se documentan los cambios y se actualiza este PRD si es necesario
\end{enumerate}

% ========================================
% 15. VALIDACION
% ========================================
\section{Criterios de Aceptación}

\begin{enumerate}
    \item Todos los endpoints funcionan correctamente según especificación
    \item Autenticación Sanctum operativa con tokens scoped
    \item Búsqueda inteligente retorna resultados relevantes en español
    \item Carrito persiste entre sesiones (Redis)
    \item Checkout crea orden y limpia carrito
    \item Notificaciones WhatsApp/SMS se envían correctamente
    \item Panel admin muestra métricas en tiempo real
    \item Imágenes se sirven correctamente vía media proxy
    \item Todos los servicios pasan health check
    \item Despliegue exitoso en VPS + Vercel
    \item HTTPS operativo en dominio de producción
    \item CORS configurado para dominio Vercel
\end{enumerate}

% ========================================
% 16. FIRMA
% ========================================
\newpage
\section{Firmas}

\vspace{2cm}

\begin{minipage}{0.45\textwidth}
\begin{center}
\rule{6cm}{0.5pt}\\[0.5cm]
\textbf{La Villa Skateboarding}\\
Cliente\\[0.3cm]
Nombre: \rule{5cm}{0.3pt}\\[0.3cm]
Fecha: \rule{5cm}{0.3pt}\\
\end{center}
\end{minipage}
\hfill
\begin{minipage}{0.45\textwidth}
\begin{center}
\rule{6cm}{0.5pt}\\[0.5cm]
\textbf{MACREAT - developer}\\
Proveedor\\[0.3cm]
Nombre: \rule{5cm}{0.3pt}\\[0.3cm]
Fecha: \rule{5cm}{0.3pt}\\
\end{center}
\end{minipage}

\vfill

\begin{center}
\textcolor{gray}{\small Documento generado el \docdate --- \apiversion}
\end{center}

\end{document}
